\documentclass{article}

\usepackage{arxiv} % same template as the iscc paper; provide arxiv.sty when compiling

\usepackage[utf8]{inputenc} % allow utf-8 input
\usepackage[T1]{fontenc}    % use 8-bit T1 fonts
\usepackage{hyperref}       % hyperlinks
\usepackage{xurl}           % simple URL typesetting
\usepackage{booktabs}       % professional-quality tables
\usepackage{amsfonts}       % blackboard math symbols
\usepackage{nicefrac}       % compact symbols for 1/2, etc.
\usepackage{microtype}      % microtypography
\usepackage{amsmath}
\usepackage{graphicx}
\usepackage{caption}
\usepackage{subcaption}
\usepackage{titling}
\usepackage{cleveref}

\graphicspath{{figures/}}

% Fallback affiliations environment (nature-style), as in the iscc paper.
\newenvironment{affiliations}{\begin{flushleft}\small}{\end{flushleft}}

\DeclareMathOperator*{\argmax}{arg\,max}
\DeclareMathOperator*{\argmin}{arg\,min}
\newcommand{\RR}{\mathbb{R}}
\newcommand{\scphytr}{\texttt{scPhyTr}}

\title{\scphytr{}: a unified phylogenetic comparative model for the evolution of
single-cell gene expression}

\author{
  Pedro F.~Ferreira$^{1,2,*}$ % TODO: add co-authors / PI before submission
}
\date{}
\date{\vspace{-5ex}}

\begin{document}
\maketitle

\begin{affiliations}
 \item $^1$Department of Biosystems Science and Engineering, ETH Zurich, Basel, Switzerland
 \item $^2$SIB Swiss Institute of Bioinformatics, Basel, Switzerland
 \item $^*$Corresponding author
\end{affiliations}

\vspace{.5cm}

\begin{abstract}
Single-cell lineage tracing now resolves the phylogeny and the transcriptome of the
same cells, raising three intertwined questions about how gene expression evolves
along a cellular genealogy: how \emph{heritable} (versus plastic) each program is,
whether it is under \emph{selection}, and which genes \emph{co-evolve}. These
questions are currently answered by separate, non-interoperable tools---phylogenetic
autocorrelation for heritability, Ornstein--Uhlenbeck model fitting for adaptation,
and local autocorrelation for gene modules---each of which treats normalized
expression as a Gaussian trait and discards the count nature of single-cell data. We
introduce \scphytr{}, a single generative model that places a latent Brownian-motion
or Ornstein--Uhlenbeck process on the lineage tree and reads it out through an
explicit count (Poisson or negative-binomial) observation model. From one fit
\scphytr{} estimates the heritability of a trait (Pagel's $\lambda$), distinguishes
neutral drift from stabilizing or adaptive selection, quantifies evolutionary rates
and the full matrix of gene--gene evolutionary correlations, and separates heritable
from within-clone plastic variation. Inference is linear time in the number of cells
and exact up to a sparsity-preserving Laplace approximation, and scales to
transcriptomes through a batched per-gene engine and phylogenetic factor analysis. On
ground-truth simulations the count model retains power to detect heritable expression
at the low sequencing depths where Gaussian-on-log methods fail; on the lineage-tracing
data from the PATH study, and transcriptome-wide on the KP-Tracer tumor atlas, \scphytr{}'s
$\lambda$ reproduces PATH's heritability ranking; and on a melanoma subline dataset it recovers adaptive expression and a heritable-versus-plastic
decomposition that flags invasion-associated genes as the most plastic. \scphytr{} is
available as an open-source Python package.
\end{abstract}

\section*{Introduction}
\label{sec:introduction}

Single-cell lineage-tracing technologies reconstruct, for thousands of cells assayed in
a single experiment, both a rooted phylogeny---from evolving CRISPR/Cas9
recorders~\cite{mckenna_gestalt_2016, raj_scgestalt_2018}, native somatic variants, or
copy-number profiles---and a paired single-cell transcriptome~\cite{wagner_lineage_2020}.
The phylogeny records how cells are related; the transcriptome records what each cell is
doing. Together they make it possible, for the first time at scale, to ask how a cell's
\emph{phenotype} evolves along its genealogy.

Three questions recur. The first is \emph{heritability versus plasticity}: does a cell
state propagate to a cell's descendants, or does it switch on timescales shorter than
cell division~\cite{schiffman_path_2024}? The second is \emph{selection}: is a gene's
expression drifting neutrally, held near an optimum by stabilizing selection, or shifted
in a particular lineage by adaptation~\cite{hirsch_emt_2025}? The third is
\emph{co-evolution}: which genes or programs change together along the tree, defining
coordinated modules of expression evolution?

Each question is, at present, addressed by a dedicated method built on a different
statistical footing. PATH quantifies heritability as the phylogenetic autocorrelation
(Moran's $I$) of a cell state and links it to a state-transition
rate~\cite{schiffman_path_2024, moran_autocorrelation_1950}. EvoGeneX and CAGEE fit
Brownian-motion and Ornstein--Uhlenbeck (OU) models---the phylogenetic comparative
methods of evolutionary biology~\cite{felsenstein_contrasts_1985, hansen_ou_1997,
butler_king_2004}---to log-transformed bulk or pseudobulk expression to call genes whose
expression has adaptively shifted between groups of
species~\cite{mendes_cagee_2023, rohlfs_eve_2015}. Hotspot scores the local phylogenetic
or spatial autocorrelation of each gene to assemble co-autocorrelated
modules~\cite{detomaso_hotspot_2021}. These tools are not interoperable: they target
different quantities, return outputs in different formats, and---critically---each
operates on a normalized, log-transformed expression value treated as a Gaussian trait,
discarding the fact that single-cell expression is observed as sparse counts whose noise
is mean-dependent. They are also predominantly univariate, so coordinated evolution is
recovered only by clustering single-gene results \emph{post hoc}, an operation that is
itself confounded by shared ancestry~\cite{felsenstein_contrasts_1985}.

The classical phylogenetic comparative framework already provides a unifying language for
all three questions: a continuous character evolving by BM or OU on a tree has a
heritability (its phylogenetic signal), a selection regime (drift versus an optimum), an
evolutionary rate, and, in the multivariate case, an evolutionary covariance between
characters. What has prevented this framework from being applied directly to single-cell
expression is threefold: the need to model sparse counts rather than a Gaussian
log-expression; the need to scale to transcriptome-wide panels of genes and to trees with
thousands of cells; and the need to estimate the multivariate structure that captures
co-evolution.

Here we introduce \scphytr{}, a phylogenetic comparative model for single-cell gene
expression that addresses these three obstacles in one coherent generative framework. A
latent BM or OU process evolves a real-valued expression state at every node of the
lineage tree, and an explicit observation model---Poisson or negative-binomial for counts,
Gaussian for directly measured traits---connects the latent state to the data at the
leaves. Because the evolutionary model and the observation model are separated, the same
inference machinery estimates, from a single fit, the heritability of a trait, its
drift-versus-selection regime, its evolutionary rate, the matrix of gene--gene
evolutionary correlations, and its within-clone plasticity---recovering, as special cases,
the quantities that PATH, EvoGeneX and Hotspot each compute separately. Inference exploits
the Markov structure of evolution on a tree to remain linear time in the number of cells,
and a batched engine and phylogenetic factor analysis extend it to transcriptome scale.

\section*{Results}
\label{sec:results}

\subsection*{One generative model for heritability, selection and co-evolution}

\scphytr{} models the log-expression of a panel of genes as a latent stochastic process
on the lineage tree, observed through a count likelihood at the leaves
(\Cref{fig:overview}). Along each branch, the latent state evolves by multivariate
Brownian motion---the null model of neutral drift, parameterized by a $p\times p$
diffusion (rate) matrix $K$---or by an Ornstein--Uhlenbeck process, which adds a pull of
strength $\alpha$ toward an optimum $\theta$ that models stabilizing selection and, with
clade-specific optima, adaptive shifts. The observed counts at each cell are a noisy
read-out of the latent state of that cell, with a per-cell size factor; cells belonging to
the same clone or subclone share that clone's latent state, so within-clone spread beyond
sampling noise is modeled as plasticity.

This single model yields, from one fit, the quantities that specialized tools compute
separately. Its diagonal $K$ entries are per-gene evolutionary rates, as in CAGEE; its
off-diagonal entries are gene--gene evolutionary correlations, the multivariate structure
that defines co-evolving programs and that univariate tools cannot access. Comparing BM,
single-optimum OU and multi-regime OU by information criteria distinguishes neutral drift
from stabilizing selection and from an adaptive shift between groups of cells, the test
that EvoGeneX performs. Fitting Pagel's $\lambda$---the proportion of trait covariance
that follows the tree---gives a heritability statistic directly analogous to PATH's
phylogenetic autocorrelation, while the OU mixing rate $\alpha$ gives the complementary
transition rate. The within-clone overdispersion gives a count-level plasticity term,
and the block structure of $K$ defines expression modules without the cell-exchangeability
null that confounds Hotspot. \scphytr{} is, in this sense, a single interpretable model
whose parameters span the union of questions addressed by the existing tools.

\subsection*{Linear-time, exact-marginal inference scales to transcriptomes}

Evolution on a tree is a Gaussian Markov random field: although the covariance between
tips is dense, the precision of the joint latent state is as sparse as the
tree~\cite{rue_held_2005}. \scphytr{} exploits this to keep all computations linear in the
number of cells. For Gaussian observations the marginal likelihood is computed exactly by
Felsenstein's pruning algorithm~\cite{felsenstein_pruning_1973}; for non-conjugate count
observations it is computed by a sparsity-preserving Laplace approximation that places a
latent at every node and never forms a dense covariance, so the cost remains $O(n)$ in the
number of cells and $O(p^3)$ in the number of genes (\Cref{sec:methods}). Parameters are
estimated by maximum marginal likelihood, either directly or by an
Expectation--Maximization algorithm whose E-step is a tree smoother and whose M-step uses
automatic differentiation; the E-step is pluggable, so the same evolutionary model can be
fit with a deterministic Laplace approximation or with asymptotically exact Markov chain
Monte Carlo. Against dense oracles, the linear-time marginal and the multivariate block
recursion agree to machine precision (to $10^{-13}$ for Gaussian observations and
$10^{-7}$ for Poisson on validation trees).

Two further components extend the framework to transcriptome scale. A batched engine
evaluates the per-gene marginal for all genes simultaneously by vectorizing the tree
recursion over the gene axis, so that the per-node computation is amortized across the
panel; in practice it fits BM, single-optimum OU and multi-regime OU and selects among
them at a few milliseconds per gene, bringing a transcriptome-wide scan from hours to
minutes. For a full $p\times p$ diffusion matrix, which has $O(p^2)$ parameters,
\scphytr{} instead fits a low-rank-plus-diagonal phylogenetic factor model
$K = WW^\top + D$ in $O(nk^3)$ time, where $k\ll p$ is the number of latent expression
programs and $D$ is a per-gene heritable idiosyncratic term.

\subsection*{Modeling counts improves heritability detection at low sequencing depth}

The distinctive methodological choice in \scphytr{}---an explicit count observation model
rather than a Gaussian model of log-normalized expression---matters most where data are
sparse. We constructed a ground-truth benchmark in which a continuous expression trait is
either heritable (Brownian on the tree) or plastic (independent across cells), generated
Poisson counts at a controlled sequencing depth, and measured each method's ability to
separate the two by the area under the receiver operating characteristic curve (AUROC),
sweeping depth from one to sixty-four mean counts per cell (\Cref{fig:depth};
\Cref{sec:methods}). All methods used the same tree and the same task; only the
observation model differed, with \scphytr{} consuming counts and the Gaussian-on-log
competitors (an EvoGeneX-style Brownian fit and PATH's Moran's $I$) consuming the
normalized trait.

At deep coverage all three methods detect heritability well (AUROC $0.93$--$0.99$ at
sixty-four counts per cell). As depth drops, the Gaussian-on-log methods degrade sharply,
falling to near-chance discrimination at one count per cell (AUROC $0.66$ for the
EvoGeneX-style fit and $0.65$ for Moran's $I$), whereas \scphytr{} retains substantial
power (AUROC $0.82$)---a gap of roughly $0.17$ AUROC attributable solely to modeling the
counts. In a complementary experiment in which we simulated a true adaptive shift in a
silenced clade and counted false-positive adaptive calls under the null, the
Gaussian-on-log test's false-positive rate inflated toward one as depth fell, while the
count model's remained stable. These results quantify, on data with known truth, the cost
of the Gaussian-on-log approximation for sparse single-cell data and the corresponding
advantage of an explicit count model.

\subsection*{\scphytr{} recovers phylogenetic heritability on lineage-tracing data}

To test whether \scphytr{} reproduces a specialized tool's read-out on that tool's own
data, we applied it to the pancreatic-cancer lineage-tracing dataset analyzed by
PATH~\cite{schiffman_path_2024, simeonov_emt_2021}, in which an epithelial-to-mesenchymal
transition (EMT) pseudotime is measured for each cell on reconstructed single-cell trees.
We fit \scphytr{}'s Pagel's $\lambda$ heritability statistic to the EMT pseudotime and
compared it, clone by clone, to PATH's Moran's $I$ (\Cref{fig:emt}).

The two statistics agree. On the clone with the most cells, \scphytr{} and PATH both call
the EMT pseudotime significantly heritable ($\lambda = 0.68$, likelihood-ratio
$p < 10^{-60}$; Moran's $I = 0.05$, permutation $p = 0.002$), and the OU mixing rate from
the same fit provides PATH's complementary transition read-out. Across clones, \scphytr{}
and PATH agree on the binary heritability call in $90\%$ of cases, and among the clones
where both statistics are statistically reliable---those with at least one hundred
cells---\scphytr{}'s $\lambda$ reproduces PATH's heritability ranking exactly (Spearman
$\rho = 1.0$). The discordances are confined to clones with fewer than fifty cells, where
the maximum-likelihood $\lambda$ overfits; this is a small-sample limitation that PATH
mitigates by bootstrapping and that \scphytr{} will adopt. \scphytr{} thus recovers PATH's
heritability measure on PATH's own data while delivering, from the same model, the
selection and rate read-outs that PATH does not provide.

To test the agreement transcriptome-wide in the sparse, in-vivo single-cell regime that
favours autocorrelation methods, we turned to the KP-Tracer mouse lung-adenocarcinoma
lineage-tracing atlas~\cite{yang_kptracer_2022}, whose reconstructed tumor trees carry
hundreds to thousands of profiled cells. On the four high-fitness tumors and the most
variable genes of each, we compared \scphytr{}'s $\lambda$ to PATH's Moran's $I$ on the
identical cell-level tree. The two agree on the binary heritability call for $87\%$ of
genes, and within each tumor $\lambda$ reproduces PATH's gene ranking (median Spearman
$\rho = 0.58$; $\rho = 0.44$--$0.77$ in three of the four tumors). In the fourth and
deepest tree both statistics still call essentially all genes heritable but their
magnitudes compress into narrow bands, so the rank correlation---though not the
calls---degrades; this exposes a property that motivates the model-based statistic, namely
that Moran's $I$ is not comparable across trees of different depth whereas $\lambda$, a
bounded proportion of variance, is. As expected for these fate-converged high-fitness
tumors, one-hot cell-state traits are dominated by a single fate and are uninformative,
and the maximum-likelihood $\lambda$ again overfits the rarest states---the same
small-sample behaviour seen on the smallest EMT clones. Because the integrated KP-Tracer
expression is distributed only in normalized form, this comparison uses \scphytr{}'s
Gaussian-trait path; the count-model advantage is the one established in simulation and on
the deeply sequenced sublines.

\subsection*{\scphytr{} detects adaptive and plastic expression in tumor subclones}

We next applied \scphytr{} to the melanoma subline dataset for which EvoGeneX was
developed~\cite{hirsch_emt_2025}, comprising single cells from clonal sublines of a mouse
melanoma model with a reconstructed subline phylogeny. Because these full-length
Smart-seq2 cells are deeply sequenced, we used the per-cell count model with the cells of
each subline treated as repeated observations of that subline's latent state, fitting BM,
single-optimum OU and adaptive two-regime OU per gene and selecting among them.

\scphytr{} jointly estimates, for every gene, an evolutionary rate, an adaptive call, and a
within-subline plasticity term. The resulting heritable-versus-plastic decomposition is
biologically coherent: the most plastic genes---those whose variation is largest within
sublines relative to the heritable component between them---are matrix and
invasion-associated genes including \emph{Thbs1}, \emph{Adamts1} and \emph{Fn1}, whereas
the most heritable genes are housekeeping genes such as \emph{Gapdh}. This recovers, within
a single count-based model, the kind of plasticity structure that motivated the original
study. We confirmed that the inferred rates are invariant to the gene-length normalization
appropriate to full-length protocols, as expected because a per-gene length factor is
absorbed by the per-gene baseline. We note one negative result that bears on
interpretation: scanning all genes detected in every subline, the genes \scphytr{} calls
adaptive are not enriched for canonical Wnt-signaling membership, the pathway highlighted
by the original analysis; the adaptive signal instead tracks expression variability. We
return to this in the Discussion.

\subsection*{\scphytr{}'s evolutionary correlation matrix defines co-evolving modules without tree confounding}

The off-diagonal of \scphytr{}'s rate matrix $K$ is a deconfounded measure of which genes
co-evolve. To benchmark it against the autocorrelation- and graph-based alternatives with a
known answer, we kept a real KP-Tracer tumor's tree and branch lengths but simulated genes
as multivariate Brownian motion with a prescribed $K$: a block of co-evolving genes plus a
majority of evolutionarily independent genes ($K$ diagonal). We then asked three methods,
at a common false-discovery threshold, to recover the gene--gene associations: \scphytr{}'s
deconfounded correlation (Felsenstein contrasts), PATH's phylogenetic cross-correlation, and
Hotspot's tree-mode local correlation.

The methods diverge sharply on the independent genes (\Cref{fig:coevo}). Shared ancestry
makes two evolutionarily independent genes co-vary at the tips, and the two methods that do
not partial it out manufacture co-evolution from it: averaged over replicates and two tumor
trees, PATH's cross-correlation falsely calls $96\%$ of independent gene pairs co-evolving
and Hotspot $73\%$, against a nominal $5\%$. \scphytr{}'s deconfounded $K$ holds the false
rate at $0.4\%$ while recovering the entire planted module (power $1.0$), separating true
from spurious co-evolution essentially perfectly (AUROC $1.0$ versus $0.74$ for PATH and
$0.5$---chance---for Hotspot). Because this comparison isolates the confounding axis, it uses
noise-free Gaussian traits, the regime most favourable to the competitors; the gap is a
property of the estimand, not of sequencing depth. The block structure of $K$ thus provides
the module definition Hotspot targets, but with the cell-exchangeability null replaced by the
phylogeny it omits.

The full $p \times p$ matrix $K$ does not scale to transcriptome-wide analysis, and a dense
estimate is in any case ill-posed: with $n$ cells the deconfounded correlation has only
$n-1$ effective degrees of freedom, so for $p > n$ it is rank-deficient and not a valid
(positive-definite) covariance. \scphytr{} therefore parameterizes $K$ in low rank through
Poisson phylogenetic factor analysis: $k$ latent factors evolve as Brownian motions on the
tree and load onto genes through a $p \times k$ matrix $W$, giving $K = W W^{\mathsf{T}}$ in
$O(N k^3)$ time and $O(p k)$ storage---the factor map, not the $p \times p$ matrix, is ever
formed. On the same real KP-Tracer trees, simulating counts from a known rank-$k$ program and
fitting the factor model recovers the planted module perfectly up to $p = 8000$ genes (block
AUROC $1.0$) in seconds, where the dense $K$ would require two orders of magnitude more
storage and is not positive-definite (\Cref{fig:pfa}). The number of factors can itself be
chosen \emph{de novo} by parallel analysis on the contrasts, which returns zero factors on
evolutionarily independent genes---where the same test on the raw leaves invents spurious
programs from shared ancestry---and recovers the true rank when a program is present.

\subsection*{\scphytr{} detects clade-specific rate shifts as fast as it estimates them, against RevBayes}

A distinct question is rate \emph{heterogeneity}: does the tempo of expression evolution
accelerate in a particular clade? The Bayesian reference for this problem is RevBayes, which
places rate shifts on the tree by reversible-jump MCMC~\cite{hohna_revbayes_2016}. \scphytr{}
fits the same multi-rate Brownian motion but localizes the shift by penalized-likelihood
forward selection, trading a posterior for speed. We compared the two on $50$-tip trees with
a single planted rate shift of varying magnitude, handing both tools the identical tree and
trait (\Cref{fig:rateshift}). At a fourfold shift both methods sat near the detection floor
(each localizing the shift to the correct clade in about a quarter of replicates), reflecting
how little a single Brownian realization on fifty tips constrains a modest rate change. At a
sixteenfold shift \scphytr{} localized the shift in every replicate against $75\%$ for
RevBayes' highest-posterior branch, and across magnitudes \scphytr{} recovered the size of
the rate change nearly without bias ($4.4\times$ and $17\times$ estimated for true $4\times$
and $16\times$), whereas RevBayes' reversible-jump posterior shrank it heavily toward the
null ($1.4\times$ and $3.6\times$). \scphytr{} did so roughly a hundred times faster (about
one second versus ninety seconds of MCMC per dataset). Its de-novo search is calibrated: with
a location penalty correcting for the choice among candidate branches, it raised no false
shift on homogeneous-rate null trees. The comparison is not all in \scphytr{}'s favour:
RevBayes solves the harder, more general problem of a full posterior over every branch's
rate and shift probability, and its shrinkage is the expected, calibrated behaviour of a
parsimonious prior, whereas \scphytr{} returns a single configuration. The point is that
\scphytr{}'s linear-time marginal makes a maximum-penalized-likelihood shift search so cheap
that it recovers the same rate-shift structure RevBayes targets, at interactive speed and
with a tighter, model-matched estimate of the rate change. On a real KP-Tracer tumor the same
trade-off recurs: for the most variable gene, \scphytr{} selects a roughly threefold rate
increase in a large subclade, whereas RevBayes' posterior places no branch above a $0.10$
shift probability---the moderate-effect regime in which, by the simulations above,
\scphytr{} is the more sensitive and RevBayes the more conservative. Without ground truth on
real data the disagreement cannot be adjudicated, but it concretely illustrates that the two
inferential philosophies can reach different calls on the same data.

\section*{Discussion}
\label{sec:discussion}

\scphytr{} brings the phylogenetic comparative methods of evolutionary biology to
single-cell gene expression in a form that is, in three respects, more general than the
specialized tools that currently partition this problem. It models the data as counts
rather than as a Gaussian function of normalized expression, which we show is the
decisive advantage at the sparse depths typical of single-cell data; it is multivariate,
estimating the gene--gene evolutionary correlations that define co-evolving programs and
that univariate tools recover only through confounded \emph{post hoc} clustering; and it
operates natively at single-cell resolution, placing a latent state at every cell of the
lineage tree. Because it is a single generative model, the heritability, selection, rate,
co-evolution and plasticity read-outs are mutually consistent and are produced by one fit,
rather than by stitching together tools with incompatible assumptions.

\scphytr{} relates to existing methods along a spectrum. PATH~\cite{schiffman_path_2024}
is a fast moment statistic---a phylogenetic autocorrelation---that measures how strongly
shared ancestry structures a trait; \scphytr{} estimates the parameters of the generating
process behind that signal, recovering PATH's heritability while adding selection, rate and
co-evolution. EvoGeneX and CAGEE~\cite{hirsch_emt_2025, mendes_cagee_2023} fit the same BM
and OU models, but to log-normalized bulk or pseudobulk expression and one gene at a time;
\scphytr{} lifts that program to single-cell-resolution, count-based and multivariate
inference. RevBayes~\cite{hohna_revbayes_2016} is the Bayesian end of this axis, fitting
rate-shift and relaxed-clock models by reversible-jump MCMC and returning full posteriors;
\scphytr{} trades that posterior for a linear-time penalized-likelihood estimate that
localizes the same shifts about seventy times faster, a difference that compounds at
single-cell tree sizes. TreeVAE and related variational approaches~\cite{lopez_treevae_2021} buy a
flexible neural decoder at the cost of an approximate variational objective and latent
variables that are not individually interpretable, whereas \scphytr{} keeps a log-linear
observation model in exchange for a near-exact marginal likelihood and identifiable
evolutionary parameters. Hotspot~\cite{detomaso_hotspot_2021} scores autocorrelation
against a cell-exchangeability null that the phylogeny violates; \scphytr{} models and
removes the tree, so its modules are defined by the evolutionary covariance rather than by
smoothed shared ancestry.

Several limitations qualify these results and define the work ahead. The count marginal is
computed by a Laplace approximation, which is exact for Gaussian observations and tight for
the log-concave Poisson likelihood but inexact in general; importance-sampling and Markov
chain Monte Carlo engines are available where higher fidelity is required. Model selection
by information criteria on approximate marginals can be liberal, and we observed that the
fraction of genes called adaptive is sensitive to the treatment of within-clone
overdispersion. The maximum-likelihood heritability statistic overfits on trees with few
tips and should be regularized by a bootstrap, as PATH does. Like the methods it
generalizes, \scphytr{} conditions on a fixed tree and, for the adaptive test, on a
supplied regime partition; propagating tree uncertainty and inferring regimes \emph{de
novo}---for example from spatial niches---are natural extensions. Finally, the negative
result on the melanoma data, in which adaptive genes were not enriched for Wnt signaling,
is a cautionary finding rather than a contradiction: \scphytr{}'s adaptive signal tracks
how variable a gene's expression is across the genealogy, which need not coincide with
membership of a curated pathway, and a faithful comparison to the original analysis will
require matching its full transcriptome-wide statistics and gene-set machinery.

The framework opens several directions. The negative-binomial batched engine will make
calibrated adaptive scans feasible transcriptome-wide; applying \scphytr{} to large,
sparse droplet-based lineage-tracing atlases is the regime where its count and single-cell
advantages should be largest, and where the comparison to PATH and Hotspot will be most
informative; and the multivariate diffusion matrix provides a principled, deconfounded
definition of expression modules that we will benchmark directly against autocorrelation-based
module detection. By expressing heritability, selection and co-evolution as parameters of a
single evolutionary model, \scphytr{} aims to make the study of how gene expression evolves
along a cellular genealogy quantitative and unified.

\section*{Methods}
\label{sec:methods}

\subsection*{The latent evolutionary model}
For a panel of $p$ genes, \scphytr{} posits a latent real-valued state $z_u\in\RR^p$ at
every node $u$ of the rooted tree $T$, interpreted as log-expression. Under Brownian
motion, the state along a branch of length $t$ takes an independent Gaussian step
$z_{\text{child}}\mid z_{\text{parent}}\sim\mathcal N(z_{\text{parent}}, t\,K)$, where the
$p\times p$ positive-definite diffusion matrix $K$ has per-gene rates on its diagonal and
evolutionary covariances off-diagonal. The Ornstein--Uhlenbeck process adds a pull toward
an optimum, $\mathrm dz = -\alpha(z-\theta)\,\mathrm ds + K^{1/2}\,\mathrm dW$, integrated
over a branch to a linear-Gaussian transition with contraction $\phi=e^{-\alpha t}$ and
variance factor $v(t)=(1-e^{-2\alpha t})/(2\alpha)$; clade-specific optima are encoded by
painting regimes on the tree~\cite{butler_king_2004}. With a fixed root the leaves are
jointly Gaussian with the separable Kronecker covariance $C\otimes K$, where $C$ is the
matrix of shared root-to-most-recent-common-ancestor times.

\subsection*{The observation model}
The observation model connects the latent leaf state to the data and factorizes across
entries given the latent. \scphytr{} implements a Poisson model $Y\sim\mathrm{Poisson}(S\,
e^{z})$ with per-cell size factors $S$; a negative-binomial extension that adds a per-gene
overdispersion to capture within-clone plasticity beyond Poisson sampling; and a Gaussian
model used for directly measured traits and as an exact oracle. Because the evolutionary
parameters live entirely in the latent model and the observation model carries none of
them, any observation type reuses the same inference engine, and the estimated $K$ is a
property of evolution rather than of measurement noise.

\subsection*{Linear-time inference}
The joint latent law is a Gaussian Markov random field whose precision has the sparsity of
the tree~\cite{rue_held_2005}. For Gaussian observations the marginal likelihood is
computed exactly by Felsenstein's pruning algorithm~\cite{felsenstein_pruning_1973,
felsenstein_contrasts_1985}. For non-conjugate observations \scphytr{} forms the
sparsity-preserving Laplace approximation: it finds the posterior mode by Newton's method
and evaluates the marginal from the tree-structured precision $Q+W$, where $W$ is the
diagonal observation curvature, by Gaussian elimination in leaves-to-root order, which
introduces no fill-in. The multivariate case generalizes this to block elimination with one
$p\times p$ Cholesky factorization per node, giving cost $O(np^3)$. Parameters are estimated
by maximum marginal likelihood, either by direct optimization or by Expectation--Maximization
whose E-step is a Rauch--Tung--Striebel tree smoother reusing the elimination factors and
whose M-step maximizes the expected complete-data log-likelihood by automatic
differentiation~\cite{dempster_em_1977, bradbury_jax_2018}. The E-step is pluggable: a
deterministic Laplace approximation, importance sampling, or Hamiltonian Monte Carlo.

\subsection*{Heritability, rate, selection and plasticity}
Heritability is quantified by Pagel's $\lambda$~\cite{pagel_lambda_1999}: the off-diagonal
entries of the tip covariance $C$ are scaled by $\lambda\in[0,1]$, with $\lambda=1$ the full
tree and $\lambda=0$ a star of independent tips, and $\lambda$ is estimated by maximum
likelihood with a likelihood-ratio test against $\lambda=0$. Because scaling only the
off-diagonals leaves a covariance that is no longer generated by a Brownian motion on the
tree, it loses the tree-Markov sparsity the pruning exploits, so $\lambda$ is the one read-out
estimated by a dense $O(n^3)$ factorization rather than the linear-time recursion---exact for
trees and subsamples up to a few thousand tips. The diagonal of $K$ gives per-gene
evolutionary rates and its off-diagonal the evolutionary correlation matrix. Selection is
assessed by fitting BM, single-optimum OU and multi-regime OU and comparing them by AIC; an
OU win indicates a pull toward an optimum and a multi-regime OU win an adaptive shift
between regimes. Clade-specific \emph{rate} heterogeneity is modelled by O'Meara's
multi-rate Brownian motion, in which the diffusion rate $\sigma^2$ changes between regimes
painted on the tree~\cite{omeara_rates_2006}; given regimes the per-regime rates are fit by
the same linear-time pruning (the root mean profiled out), and when the shifts are unknown
they are found \emph{de novo} by greedy forward selection that adds the branch most
improving the BIC, with a $2\log m$ location penalty per shift ($m$ candidate branches) that
corrects for selecting the best of many branches and controls the null false-positive rate;
this is the maximum-penalized-likelihood counterpart to RevBayes' reversible-jump MCMC over
the number and location of shifts~\cite{hohna_revbayes_2016, khabbazian_l1ou_2016}. Within-clone
plasticity is the negative-binomial overdispersion estimated
jointly with the latent model, and the heritable-versus-plastic decomposition of a gene is
the ratio of its between-clone heritable variance to its within-clone variance.

\subsection*{Scalability}
The batched engine evaluates the per-gene Laplace marginal for all genes at once by carrying
a gene axis through the tree recursion, so each per-node scalar operation becomes a vector
operation over genes and the Python-level traversal is paid once for the whole panel; it
fits BM/OU/OU-regime per gene by maximizing the marginal over a grid of $(\alpha,\sigma^2)$
with the optimum profiled in closed form. The phylogenetic factor model constrains
$K = WW^\top + D$ with $k\ll p$ shared latent programs and a per-gene heritable idiosyncratic
term $D$, fit by Laplace-EM in $O(nk^3)$.

\subsection*{Benchmark design}
The heritability-detection benchmark simulates, on a fixed tree, traits that are either
heritable (a Brownian draw with covariance $C$) or plastic (independent across cells),
generates Poisson counts at a series of mean depths, and scores each method by the AUROC of
separating heritable from plastic traits. \scphytr{} scores a trait by its Pagel's $\lambda$;
the EvoGeneX surrogate scores it by the Brownian-motion fit to the per-subline mean of
$\log(1+\text{count})$; PATH scores it by Moran's $I$ with a node-depth phylogenetic weight
matrix and a permutation test. The adaptive-detection experiment simulates a down-shift of
expression in a chosen clade under the alternative and no shift under the null, and reports
power and false-positive rate as a function of depth.

\subsection*{Datasets}
The melanoma subline data are the full-length Smart-seq2 single-cell expression of clonal
sublines of a mouse melanoma model~\cite{hirsch_emt_2025}; raw RSEM expected counts were
obtained from Gene Expression Omnibus accession GSE215960, mapped to subline labels, and
analyzed on the published consensus subline phylogeny. The EMT lineage-tracing data are the
pancreatic-cancer single-cell trees and EMT pseudotime from the PATH
study~\cite{schiffman_path_2024, simeonov_emt_2021}; per-clone trees and per-cell pseudotime
were exported from the published data objects. External datasets are not redistributed with
the software.

\subsection*{Comparison methods}
PATH was reimplemented as Moran's $I$ with a shared-ancestry phylogenetic weight matrix and a
permutation null. The EvoGeneX comparison uses \scphytr{}'s own Gaussian phylogenetic
BM/OU machinery applied to the per-subline mean of $\log(1+\text{count})$, a faithful and
favorable surrogate for the Gaussian-on-log approach that isolates the effect of the
observation model. Concordance on real data is reported as the agreement of binary
heritability calls and the Spearman correlation of the heritability statistics across clones.
Hotspot~\cite{detomaso_hotspot_2021} was run in its tree (PhyloVision) mode through the
\texttt{hotspotsc} package. RevBayes~\cite{hohna_revbayes_2016} (v1.4.0) was run on the
identical tree and trait as a random local clock: each branch carries a reversible-jump
mixture rate multiplier (no shift, or a lognormal draw) that is inherited by its descendants,
so a single jump at a clade's base raises the whole clade's rate; the per-branch posterior
shift probabilities and rates summarize a $5\times10^4$-generation chain. Because RevBayes'
Brownian REML model requires a bifurcating tree, multifurcating lineage trees were resolved
to binary with floored branch lengths before both tools were run on real data.

\subsection*{Implementation and code availability}
\scphytr{} is implemented in Python on top of NumPy, SciPy and JAX, and is available as the
\texttt{scphytr} package. % TODO: public repository URL and DOI before submission

\section*{Competing interests}
The authors declare that they have no competing interests.

\section*{Acknowledgements}
% TODO: funding and acknowledgements.

\bibliographystyle{unsrt}
\bibliography{references}

% ---------------------------------------------------------------------------
% Figures (to be generated from analysis/ scripts before submission)
% ---------------------------------------------------------------------------
\begin{figure}[t]
  \centering
  % \includegraphics[width=\textwidth]{fig1_overview.pdf}
  \caption{\textbf{\scphytr{} models single-cell gene expression as a latent evolutionary
  process on the lineage tree.} A latent Brownian-motion or Ornstein--Uhlenbeck state
  evolves along the tree and is observed as counts at the leaves; one fit yields
  heritability (Pagel's $\lambda$), selection (drift vs.\ OU optimum), evolutionary rate and
  gene--gene correlations ($K$), and within-clone plasticity, spanning the read-outs of
  PATH, EvoGeneX/CAGEE and Hotspot. Source: schematic.}
  \label{fig:overview}
\end{figure}

\begin{figure}[t]
  \centering
  % \includegraphics[width=0.7\textwidth]{fig2_depth_auroc.pdf}
  \caption{\textbf{The count model retains heritability-detection power at low sequencing
  depth.} AUROC for separating heritable from plastic traits versus mean counts per cell, on
  ground-truth simulations, for \scphytr{} (counts), an EvoGeneX-style Gaussian-on-log fit,
  and PATH (Moran's $I$). At one count per cell \scphytr{} reaches AUROC $0.82$ versus
  $0.66$ and $0.65$ for the Gaussian-on-log methods; all converge by sixty-four counts per
  cell. Source: \texttt{analysis/benchmark/sim\_heritability.py}.}
  \label{fig:depth}
\end{figure}

\begin{figure}[t]
  \centering
  % \includegraphics[width=0.7\textwidth]{fig3_emt_concordance.pdf}
  \caption{\textbf{\scphytr{}'s heritability statistic reproduces PATH's on the EMT
  lineage-tracing data.} Per-clone Pagel's $\lambda$ (\scphytr{}) versus Moran's $I$ (PATH)
  for the EMT pseudotime; the two agree on the heritability call in $90\%$ of clones and rank
  reliable clones identically (Spearman $\rho = 1.0$ for clones with $\geq 100$ cells).
  Source: \texttt{analysis/benchmark/emt\_concordance.py}.}
  \label{fig:emt}
\end{figure}

\begin{figure}[t]
  \centering
  % \includegraphics[width=\textwidth]{fig4_melanoma.pdf}
  \caption{\textbf{Heritable-versus-plastic decomposition of melanoma subline expression.}
  For each gene, heritable (between-subline) versus plastic (within-subline) variance from a
  single count-based \scphytr{} fit; invasion-associated genes (\emph{Thbs1}, \emph{Adamts1},
  \emph{Fn1}) are the most plastic and housekeeping genes the most heritable. Source:
  \texttt{analysis/melanoma/heritable\_vs\_plastic.py}.}
  \label{fig:melanoma}
\end{figure}

\begin{figure}[t]
  \centering
  % \includegraphics[width=0.7\textwidth]{fig5_coevolution.pdf}
  \caption{\textbf{\scphytr{}'s evolutionary correlation matrix recovers co-evolving modules
  without tree confounding.} Genes simulated as multivariate Brownian motion with a known
  co-evolving block on real KP-Tracer tumor trees. Power to recover the block versus the
  false-positive rate among evolutionarily independent gene pairs (nominal $5\%$), at a common
  BH-FDR. \scphytr{}'s deconfounded $K$ (Felsenstein contrasts) keeps the false rate at
  $0.4\%$ at full power (AUROC $1.0$); PATH's phylogenetic cross-correlation ($96\%$ false,
  AUROC $0.74$) and Hotspot's tree-mode local correlation ($73\%$ false, AUROC $0.5$) invent
  co-evolution from shared ancestry. Source:
  \texttt{analysis/benchmark/kptracer\_coevolution.py}.}
  \label{fig:coevo}
\end{figure}

\begin{figure}[t]
  \centering
  % \includegraphics[width=0.85\textwidth]{fig6_pfa_scaling.pdf}
  \caption{\textbf{Low-rank factorization scales the evolutionary correlation matrix to the
  transcriptome.} Counts simulated from a known rank-$k$ program on real KP-Tracer trees
  ($n = 400$ cells). \textbf{(A)} Poisson phylogenetic factor analysis recovers the planted
  module (block AUROC) up to $p = 8000$ genes in seconds; the dense $p \times p$ correlation
  matches recovery only where it can be formed and is rank-deficient (not positive-definite)
  for $p > n$. \textbf{(B)} Stored parameters: the factor model grows as $p\,k$ while the dense
  $K$ grows as $p^2/2$ (two orders of magnitude more at $p = 4000$). Source:
  \texttt{analysis/benchmark/kptracer\_coevolution.py} (\texttt{pfa}).}
  \label{fig:pfa}
\end{figure}

\begin{figure}[t]
  \centering
  % \includegraphics[width=\textwidth]{fig7_rateshift.pdf}
  \caption{\textbf{Clade-specific rate-shift detection: \scphytr{} versus RevBayes.} A
  continuous trait simulated with one rate shift of known magnitude on $50$-tip trees, given
  identically to both tools. \textbf{(A)} Fraction of replicates localizing the shift to the
  true clade (\scphytr{}: BIC-selected shift; RevBayes: branch posterior shift probability
  $>0.5$) versus shift magnitude. \textbf{(B)} Estimated versus true clade/background rate
  ratio: \scphytr{} (penalized likelihood) tracks the truth, RevBayes' reversible-jump
  posterior shrinks toward the null. \scphytr{} runs in $\sim$1\,s per dataset against
  $\sim$85\,s of MCMC. Source: \texttt{analysis/benchmark/revbayes\_rate\_shifts.py}.}
  \label{fig:rateshift}
\end{figure}

\end{document}
